% TOP_PROBLEM: {"id": 10, "title": "The Hadwiger-Nelson Problem", "category_id": 2, "category": {"id": 2, "name": "combinatorics", "display_name": "Combinatorics", "description": "Counting problems, graph theory, discrete structures.", "slug": "combinatorics", "order_index": 2, "created_at": "2026-07-31T15:26:25.671Z"}, "status": "open"}
% FIRST_POSTED: 2026-09-25
% ENTRY_KIND: statement_only
% !TeX program = lualatex
% Definition and sources only; this file is not an LLM solution attempt.
\documentclass[11pt]{article}
\usepackage[margin=25mm]{geometry}
\usepackage{fontspec}
\setmainfont{Latin Modern Roman}
\usepackage{amsmath,amssymb,mathtools}
\usepackage{unicode-math}
\setmathfont{Latin Modern Math}
\usepackage{xurl,hyperref}
\hypersetup{colorlinks=true,urlcolor=blue}
\setcounter{secnumdepth}{0}
\setlength{\parindent}{0pt}
\setlength{\parskip}{5pt}
\setlength{\emergencystretch}{3em}
\begin{document}
\section*{\texorpdfstring{The Hadwiger-Nelson Problem}{The Hadwiger-Nelson Problem}}\label{10}
\subsection{Definitions and mathematical statement}
The unit-distance graph of the plane has vertex set ${\mathbb{R}}^2$, with an edge between $x,y$ precisely when $|x-y|=1$. Its chromatic number is
\[
 \chi({\mathbb{R}}^2)=\min\{k\in{\mathbb{N}}:\exists c:{\mathbb{R}}^2\to\{1,\ldots,k\},
 \ |x-y|=1\Rightarrow c(x)\ne c(y)\}.
\]
Determine this integer in ordinary ZFC. The coloring is an arbitrary function; its color classes need not be measurable, open, or periodic. Colorings of only a finite set of points provide obstructions but are not by themselves colorings of the whole plane. The exact-value target is separate from merely improving a lower or upper bound.
\par\smallskip\textit{Formulation and concept references:} \href{https://arxiv.org/abs/1804.02385}{[S1]}.

\subsection{Short English statement}
What is the smallest number of colors that can color the entire plane without giving the same color to two points exactly one unit apart?
\subsection{Sources}
\begin{itemize}
\item[S1]de Grey, The chromatic number of the plane is at least 5 (2018).
\url{https://arxiv.org/abs/1804.02385}.
\textit{Formulation source.}
\item[S2]Status source for Chromatic Number of the Plane (Hadwiger--Nelson Problem).
\url{https://arxiv.org/abs/2108.12861}.
\textit{Further reference; not a proof endorsement.}
\end{itemize}
\end{document}
